% Chapter 12: Automation & Pipelines
\chapter{Automation and Pipelines}
\label{ch:automation}

\begin{itshape}
Your team merges 20 PRs per week. Each needs type checking, test runs,
and documentation updates. Manual review is not scaling.
Claude can do this work---but not through interactive sessions.
\end{itshape}

Claude Code is not limited to interactive use.
Its headless mode, fan-out capabilities, and the Claude Agent SDK
enable production automation pipelines that run without human interaction.
The extension mechanisms from \cref{ch:extending}---especially hooks---provide
the building blocks. This chapter assembles them into pipelines.

% -------------------------------------------------------------------
\section{Choosing Your Automation Mode}
\label{sec:automation-strategy}

\marginnote[Tip]{Match the automation mode to the
failure cost. Higher stakes = more control.}

Claude Code offers five automation modes, each with different tradeoffs.
Choosing the wrong one either over-engineers simple tasks or
under-controls risky ones.

\begin{decisionbox}[Which Automation Mode?]
\textbf{Need a quick one-off script?}\\
$\rightarrow$ \textbf{Headless mode} (\code{claude -p}). Fastest to set up.
Runs a single prompt, returns output. Add \code{--bare} for
speed, \code{--json-schema} for structured output.

\textbf{Need to process many files with the same operation?}\\
$\rightarrow$ \textbf{Fan-out pattern}. Shell loop + headless mode.
You define the units; Claude processes each independently.
Test on 2--3 files before scaling.

\textbf{Need Claude to figure out what needs changing?}\\
$\rightarrow$ \textbf{\code{/batch}}. Claude researches the codebase,
decomposes the work, spawns parallel agents.
You review the plan before execution.

\textbf{Need programmatic control, custom tools, or error handling?}\\
$\rightarrow$ \textbf{Agent SDK}. Full Python/TypeScript API.
Custom tool definitions, hook callbacks, session management.
Production-grade.

\textbf{Need it to run on a schedule without your machine?}\\
$\rightarrow$ \textbf{Cloud scheduled tasks} or \textbf{GitHub Actions}.
Anthropic infrastructure runs the prompt on a cron schedule.
Each run clones fresh.
\end{decisionbox}

\subsection*{The Progression}

Most automation starts as a one-off \code{claude -p} call that you
run manually. When it works, it migrates:

\begin{enumerate}
  \item \textbf{Manual headless}: \code{claude -p "add type hints to auth.py"} ---
    you run it, review output, iterate on prompt.
  \item \textbf{Scripted fan-out}: wrap in a shell loop, add
    \code{--allowedTools} for safety, \code{--output-format json}
    for pipeline consumption.
  \item \textbf{CI integration}: add to GitHub Actions, trigger on PR.
    Now it runs without you.
  \item \textbf{Scheduled}: promote to cloud scheduled task for
    recurring work (daily reviews, weekly audits).
  \item \textbf{Agent SDK}: when the shell script gets complex enough
    to need conditionals, retries, or custom tools, rewrite in Python.
\end{enumerate}

\noindent Each step adds reliability and removes manual intervention.
The mistake is jumping to step 5 before proving the concept at step 1.

\subsection*{Trust Gates Between Stages}

\marginnote[Tip]{Do not promote a pipeline stage
until the current stage has succeeded
N times without intervention.}

Each promotion in the progression should be gated on evidence
that the current stage works reliably:

\begin{description}
  \item[Manual $\rightarrow$ Scripted] Gate: the prompt has produced
    correct results on 5+ files without modification. If you are still
    editing the prompt after each run, it is not ready to script.

  \item[Scripted $\rightarrow$ CI] Gate: the script has run
    end-to-end 3+ times with no manual intervention. Output validation
    is automated (not ``I eyeballed the diff''). Tool scoping
    (\code{--allowedTools}) is explicit.

  \item[CI $\rightarrow$ Scheduled] Gate: CI runs have produced
    correct results for 2+ weeks. Failure handling is tested
    (you know what happens when Claude hits a rate limit or an
    unexpected file). Outputs are PRs, not direct commits.

  \item[Any stage $\rightarrow$ Agent SDK] Gate: the bash script
    has grown to include conditionals, error handling, or retry logic
    that bash handles poorly. The prompt is stable. If you are still
    iterating on the prompt, stay in bash---it is faster to edit.
\end{description}

\practitioner{The trust gate pattern prevents the most common
automation failure: a prompt that works 90\% of the time deployed
to CI where it fails silently 10\% of the time. The 10\% failure
rate is invisible in manual runs (you catch and fix it) but
compounds in automation (10 failures per 100 PRs, each requiring
manual cleanup).}

\subsection*{Handling Failures in Unattended Pipelines}

\practitioner{Unattended pipelines need failure modes designed in advance.
Claude cannot ask you for help when you are not watching.}

\begin{itemize}
  \item \textbf{Bound execution}: Use \code{--max-turns} and
    \code{--max-budget-usd} to prevent runaway sessions.
    A stuck agent consuming tokens for 30 minutes is expensive.
  \item \textbf{Scope tools}: \code{--allowedTools Read,Edit} prevents
    an agent from running arbitrary commands. Match tools to
    the task, not to convenience.
  \item \textbf{Validate output}: Parse \code{--output-format json}
    results programmatically. If the output does not match expected
    structure, fail loudly---do not commit partial results.
  \item \textbf{Gate before merge}: CI pipelines should produce PRs,
    not direct commits. A human reviews before merge.
    The automation eliminates toil, not judgment.
\end{itemize}

% -------------------------------------------------------------------
\section{Headless Mode}
\label{sec:headless}

\marginnote[Official]{\code{claude -p "prompt"} for non-interactive use.
\code{--output-format json|stream-json|text}.\sourceurl{https://code.claude.com/docs/en/overview}}

\official{Use \code{claude -p} for non-interactive automation.\sourceurl{https://code.claude.com/docs/en/overview}}
\index{headless mode}

\begin{description}
  \item[\code{-p "prompt"}] Single prompt, no interaction.
  \item[\code{--output-format json}] Structured output for pipeline consumption.
  \item[\code{--allowedTools}] Restrict to specific tools for safety.
  \item[\code{--model}] Override model per invocation.
\end{description}

\subsection{Fan-Out Pattern}
\index{headless mode!fan-out}

\begin{fullwidth}
\begin{beforeafter}
\before
\begin{minted}{text}
Manual: open Claude, paste file, ask for type hints,
copy result, repeat 47 times.
\end{minted}

\after
\begin{minted}{bash}
# Test on 3 files first
for f in $(find src/ -name "*.py" | head -3); do
  claude -p "Add type hints to $f" \
    --allowedTools Read,Edit
done

# Review results, then scale
find src/ -name "*.py" | while read f; do
  claude -p "Add type hints to $f" \
    --allowedTools Read,Edit \
    --output-format json >> results.jsonl
done
\end{minted}
\end{beforeafter}
\end{fullwidth}

\official{The \code{--allowedTools} flag scopes what Claude can do in unattended
mode.\sourceurl{https://code.claude.com/docs/en/best-practices}}
\index{headless mode!allowedTools@\texttt{--allowedTools}}
Use glob patterns for fine-grained control:

\begin{minted}{bash}
# Allow Read and Edit only (safest for code modifications)
--allowedTools "Read,Edit"

# Allow Bash but only for specific commands
--allowedTools 'Read,Edit,Bash(git commit *)'
\end{minted}

\begin{warnbox}[Test Before Scaling]
Always test your fan-out prompt on 2--3 files before running at scale.
A prompt that works on \code{src/utils.py} may produce unexpected results
on \code{src/models/network.py}. Review the first few outputs manually,
then scale with confidence.
\end{warnbox}

\margintip{Combine \code{--allowedTools} with \code{--output-format json}
for pipeline-friendly output. Parse results with \code{jq} for automated
quality checks before committing.}

\subsection{Permissions in Headless Mode}

\begin{warnbox}[Headless Permissions]
In headless mode, Claude still respects permission rules.
For trusted CI/CD environments, use \code{--dangerously-skip-permissions}
to bypass interactive prompts.
This flag should \textbf{only} be used in environments you fully control
(CI runners, Docker containers). Never expose it in user-facing scripts.
\end{warnbox}

\subsection{Bare Mode for Scripted Calls}
\label{subsec:bare-mode}

\official{The \code{--bare} flag skips auto-discovery of hooks, skills,
plugins, MCP servers, auto memory, and CLAUDE.md so scripted calls
start faster.\sourceurl{https://code.claude.com/docs/en/cli-reference}}
\index{bare mode}\index{headless mode!bare mode}

In bare mode, Claude has access to only Bash, file read, and file edit
tools. This makes it ideal for CI/CD pipelines, SDK-style automation,
and \code{-p} calls where startup speed matters more than full context.

\begin{minted}{bash}
# ~14% faster to API request than a standard session
claude --bare -p "explain this function" < src/main.py
\end{minted}

\margintip{Use \code{--bare} when you don't need hooks, skills, or MCP.
It sets \code{CLAUDE\_CODE\_SIMPLE} internally.}

\practitioner{Bare mode pairs naturally with \code{--json-schema} for
structured output and \code{--max-turns} for bounded execution---the
three flags together define a predictable, fast automation primitive.}

% -------------------------------------------------------------------
\section{Claude Agent SDK}
\label{sec:agent-sdk}

\marginnote[Official]{Same agent loop as Claude Code, programmable in
Python and TypeScript.\sourceurl{https://platform.claude.com/docs/en/agent-sdk/overview}}

\begin{decisionbox}[Headless Mode vs.\ Agent SDK]
\textbf{Headless mode} (\code{claude -p}): quick automation, scripting,
CI/CD integration. Uses Claude Code's existing configuration.
Best for: fan-out, batch operations, simple pipelines.

\textbf{Agent SDK}: programmatic control, custom tools, production deployment.
Best for: custom applications, complex orchestration, SaaS integration.

\textbf{Rule}: Start with headless mode. Move to Agent SDK when you need
custom tool definitions or programmatic control flow.
\end{decisionbox}

\official{The Agent SDK provides the same tools, agent loop, and context
management that power Claude Code, programmable in Python and
TypeScript.\sourceurl{https://platform.claude.com/docs/en/agent-sdk/overview}}
\index{Agent SDK}

\begin{minted}{bash}
# Install
pip install claude-agent-sdk          # Python
npm install @anthropic-ai/claude-agent-sdk  # TypeScript
\end{minted}

\begin{fullwidth}
\begin{beforeafter}[Headless Mode vs.\ Agent SDK]
\before
\begin{minted}{bash}
# Headless: shell scripting, quick automation
claude -p "Add type hints to $f" \
  --allowedTools Read,Edit
\end{minted}

\after
\begin{minted}{python}
# Agent SDK: programmatic control
from claude_agent_sdk import query, ClaudeAgentOptions
import asyncio

async def main():
    async for msg in query(
        prompt="Add type hints to auth.py",
        options=ClaudeAgentOptions(
            allowed_tools=["Read", "Edit"],
        ),
    ):
        if hasattr(msg, "result"):
            print(msg.result)

asyncio.run(main())
\end{minted}
\end{beforeafter}
\end{fullwidth}

Key capabilities:
\begin{description}
  \item[Built-in tools] Read, Write, Edit, Bash, Glob, Grep, WebSearch, WebFetch---no
    implementation needed.
  \item[Hooks] Callback functions for PreToolUse, PostToolUse, Stop, etc.---validate,
    log, or block agent behavior programmatically.
  \item[Subagents] Spawn specialized agents with their own context and tools.
  \item[MCP] Connect to external services (databases, browsers, APIs).
  \item[Sessions] Resume or fork sessions to maintain context across exchanges.
  \item[Filesystem config] CLAUDE.md, skills, and memory all work via
    \code{setting\_sources=["project"]}.
\end{description}

\margintip{Start with headless mode (\code{claude -p}). Move to the Agent SDK
when you need custom tool definitions, programmatic hooks, or production
error handling.}

% -------------------------------------------------------------------
\section{Pipeline Design Patterns}
\label{sec:pipeline-patterns}

\subsection{Traffic-Light Dashboards}

\practitioner{Replace pass/fail with traffic-light thresholds (GREEN/YELLOW/RED).}
This provides nuance: 78\% coverage is not ``failing''---it is
YELLOW, meaning ``acceptable but track for improvement.''

\begin{minted}{text}
=== Project Health Dashboard ===
Test Coverage:    82% [GREEN]
Lint Warnings:     3  [YELLOW]
Documentation:   91% [GREEN]
Model Metrics:   AUC 0.84 [GREEN]
Overall:         HEALTHY (3 GREEN, 1 YELLOW)
\end{minted}

\subsection{Pipeline Directionality}

\begin{warnbox}[Pipeline Direction]
Never re-run upstream steps after editing downstream artifacts.
Mark generated files with \code{\# GENERATED --- do not edit manually}.
This is anti-pattern \#5 from \cref{ch:antipatterns}.
\end{warnbox}

% -------------------------------------------------------------------
\section{Scheduled Tasks}
\label{sec:scheduled-tasks}

\marginnote[Official]{Three tiers of recurring automation, each with different
tradeoffs.\sourceurl{https://code.claude.com/docs/en/web-scheduled-tasks}}

\official{Claude Code offers three ways to schedule recurring
work.\sourceurl{https://code.claude.com/docs/en/web-scheduled-tasks}}

\begin{fullwidth}
\begin{tabular}{@{}lp{2.5cm}p{2cm}p{2cm}p{3.5cm}@{}}
\toprule
& \textbf{Cloud} & \textbf{Desktop} & \textbf{\code{/loop}} & \textbf{Best For} \\
\midrule
Runs on & Anthropic cloud & Your machine & Your machine & \\
Requires machine & No & Yes & Yes & \\
Persistent & Yes & Yes & No (session) & \\
Local files & No (fresh clone) & Yes & Yes & \\
Min interval & 1 hour & 1 minute & 1 minute & \\
\cmidrule{5-5}
& & & & Cloud: daily PR reviews, dependency audits \\
& & & & Desktop: local builds, file watches \\
& & & & \code{/loop}: deploy polling, session-scoped checks \\
\bottomrule
\end{tabular}
\end{fullwidth}

\begin{description}
  \item[Cloud (\code{/schedule})] Create from CLI, web, or Desktop app.\index{scheduled tasks!Cloud}
    Runs on Anthropic infrastructure---works even when your machine is off.
    Each run clones the repo fresh, executes the prompt, and creates a session
    you can review.
  \item[Desktop] Schedule tasks that run on your machine with access to
    local files and MCP servers.\index{scheduled tasks!Desktop}
  \item[\code{/loop}] Bundled skill.\index{scheduled tasks!loop@\texttt{/loop}} Runs a prompt repeatedly on an interval
    while the session stays open: \code{/loop 5m check if the deploy finished}.
    Session-scoped---stops when you close the session.
\end{description}

\margintip{Use \textbf{Cloud} for work that must happen reliably (daily reviews,
weekly audits). Use \textbf{\code{/loop}} for quick polling during active work.
Use \textbf{Desktop} when you need local filesystem access on a schedule.}

% -------------------------------------------------------------------
\section{GitHub Actions Integration}
\label{sec:github-actions}

\marginnote[Official]{Run \code{/install-github-app} to set up Claude Code
as a GitHub Action.\sourceurl{https://code.claude.com/docs/en/github-actions}}

\official{Claude Code GitHub Actions brings AI-powered automation to your
GitHub workflow.\sourceurl{https://code.claude.com/docs/en/github-actions}}
\index{GitHub Actions}
Mention \code{@claude} in any PR or issue comment, and Claude can analyze
code, create pull requests, implement features, and fix bugs---following
your CLAUDE.md guidelines.

\begin{minted}{yaml}
# .github/workflows/claude.yml
name: Claude Code
on:
  issue_comment:
    types: [created]
  pull_request_review_comment:
    types: [created]
jobs:
  claude:
    runs-on: ubuntu-latest
    steps:
      - uses: anthropics/claude-code-action@v1
        with:
          anthropic_api_key: ${{ secrets.ANTHROPIC_API_KEY }}
\end{minted}

\begin{description}
  \item[Setup] \code{/install-github-app} walks through installation.
    Alternatively, install the Claude GitHub app manually and add
    \code{ANTHROPIC\_API\_KEY} as a repository secret.
  \item[Customization] Use \code{claude\_args} to pass CLI flags:
    \code{--max-turns}, \code{--model}, \code{--allowedTools}.
  \item[Cloud providers] Supports direct API, AWS Bedrock, and
    Google Vertex AI authentication.
\end{description}

% -------------------------------------------------------------------
\section{Remote Control and Web Sessions}
\label{sec:remote-control}

\official{Remote Control connects claude.ai/code or the Claude mobile app
to a session running on your machine.\sourceurl{https://code.claude.com/docs/en/remote-control}}
\index{remote control}
Start a task at your desk, then pick it up from your phone or another browser.
Your local filesystem, MCP servers, and project configuration stay available.

\begin{description}
  \item[\code{/remote-control}] Enable from an existing session. Also
    available as \code{claude --remote-control} or server mode
    \code{claude remote-control}.
  \item[Claude Code on the web] Runs on Anthropic cloud infrastructure---no
    local machine needed. Use for repos you don't have cloned or for
    parallel cloud tasks.
\end{description}

\marginnote[Cross-Ref]{For the full comparison of Remote Control, Dispatch,
Channels, Slack, and Scheduled Tasks, see the official
docs.\sourceurl{https://code.claude.com/docs/en/remote-control}}

% -------------------------------------------------------------------
\begin{keyinsight}
Traditional automation breaks on unexpected input---a file with an
unusual encoding, a function signature it was not designed for.
Claude-powered automation bends: it reads the unusual file, adapts
its approach, and continues. This means the boundary between
``common case'' (automated) and ``exception'' (manual) shifts
dramatically. Tasks that required human judgment in a bash pipeline---
like deciding how to refactor a function with complex dependencies---
become part of the automated path. Design your pipelines with this
wider boundary in mind: the trust gates (\cref{sec:automation-strategy})
are the new exception handling.
\end{keyinsight}

% -------------------------------------------------------------------
\begin{trybox}
\begin{enumerate}
  \item Write a \code{claude -p} command that adds type hints to one file in your project.
  \item Test it. Does the output compile? Do tests pass?
  \item If yes: scale to 3 files. Review the results.
  \item If no: refine the prompt. What did you need to specify that you assumed?
\end{enumerate}
\end{trybox}
